% Part: history
% Chapter: biographies
% Section: alonzo-church

\documentclass[../../../include/open-logic-section]{subfiles}

\begin{document}

\olfileid{his}{bio}{chu}

\olsection{অ্যালোনজো চার্চ}

\olphoto{church-alonzo}{অ্যালোনজো চার্চ}

অ্যালোনজো চার্চের জন্ম ১৯০৩ সালের ১৪ জুন ওয়াশিংটন,
ডিসি-তে। ছোটবেলায় এয়ারগান-সংক্রান্ত একটি দুর্ঘটনায়
তাঁর এক চোখের দৃষ্টিশক্তি নষ্ট হয়। ১৯২০ সালে কানেকটিকাটে
প্রস্তুতিমূলক বিদ্যালয়ের পড়াশোনা শেষ করে সেই বছরই তিনি
প্রিন্সটনে বিশ্ববিদ্যালয়ের পড়া শুরু করেন। ১৯২৭ সালে
তাঁর ডক্টরেট সম্পূর্ণ হয়। কয়েক বছর বিদেশে কাটিয়ে চার্চ
প্রিন্সটনে ফিরে আসেন। তিনি অত্যন্ত ভদ্র ও যত্নশীল মানুষ
হিসেবে পরিচিত ছিলেন। ব্ল্যাকবোর্ডে তাঁর লেখা ছিল
নিখুঁত; গুরুত্বপূর্ণ কাগজপত্র সংরক্ষণের জন্য তিনি
সেগুলির উপর যত্ন করে ডুকো সিমেন্টের (এক ধরনের স্বচ্ছ
আঠা) প্রলেপ দিতেন। প্রাতিষ্ঠানিক কাজের বাইরে তিনি
কল্পবিজ্ঞান পত্রিকা পড়তে ভালোবাসতেন; লেখায় ভুল
চোখে পড়লে সম্পাদকদের চিঠি লিখতেও দ্বিধা করতেন না।

চার্চের প্রাতিষ্ঠানিক অবদান ছিল অসামান্য। তাঁর ছাত্র
স্টিফেন ক্লিনি ও বার্কলি রসারের সঙ্গে তিনি কার্যকর
গণনসাধ্যতার একটি তত্ত্ব, ল্যাম্বডা ক্যালকুলাস, নির্মাণ
করেন; অ্যালান টুরিংয়ের টুরিং যন্ত্র নির্মাণ থেকে এই কাজ
স্বতন্ত্র ছিল। গণনসাধ্যতার এই দুই সংজ্ঞা সমতুল্য।
এগুলি থেকে বর্তমানে \emph{চার্চ--টুরিং থিসিস} নামে
পরিচিত ধারণাটির জন্ম: স্বাভাবিক সংখ্যার একটি অপেক্ষক
ঠিক তখনই কার্যকরভাবে গণনসাধ্য, যখন সেটি টুরিং যন্ত্র
(অথবা ল্যাম্বডা ক্যালকুলাস) দিয়ে গণনা করা যায়। তিনি
\emph{চার্চের উপপাদ্য} নামেও পরিচিত একটি ফল প্রমাণ
করেন: প্রথম-ক্রমের সূত্রের বৈধতা নির্ণয়ের সিদ্ধান্ত-সমস্যা
অসমাধানযোগ্য।

বার্ধক্যেও চার্চ তাঁর কাজ চালিয়ে যান। ১৯৬৭ সালে তিনি
প্রিন্সটন ছেড়ে ইউসিএলএ-তে যান এবং ১৯৯০ সালে অবসর
নেওয়া পর্যন্ত সেখানে অধ্যাপক ছিলেন। ১৯৯৫ সালের
১ আগস্ট ৯২ বছর বয়সে তাঁর মৃত্যু হয়।

\begin{reading}
চার্চের সংক্ষিপ্ত জীবনী জানতে \citet{EndertonND} দেখো।
ল্যাম্বডা ক্যালকুলাস এবং সিদ্ধান্ত-সমস্যা (চার্চের থিসিস)
নিয়ে তাঁর মূল রচনাগুলি হলো \citet{Church1936,Church1936a}।
১৯৩০-এর দশকে প্রিন্সটনের গণিতচর্চার পরিবেশ নিয়ে
চার্চের একটি সাক্ষাৎকার \citet{Aspray1984}-এ নথিবদ্ধ
আছে। ১৯৩৬ থেকে ১৯৭৯ সাল পর্যন্ত চার্চ \emph{Journal of
Symbolic Logic}-এর বইসমালোচনার একটি ধারাবাহিক রচনা
লিখেছিলেন। জন ম্যাকফারলেনের ওয়েবসাইটে সেগুলির
সবই সংরক্ষিত আছে \citep{MacFarlane2015}।
\end{reading}

\end{document}
